\documentstyle{article}
\pagestyle{empty}
\begin{document}
\section*{
Computation of coefficients of the normal form in a singular case
}
\centerline{\bf
S. Yu. Sadov,}
\vskip 2ex
%
\centerline{\bf
Institute of Applied Mathematics,}
\centerline{\bf
Moscow, Russia
}
\vskip 3ex
%
Satellite motion about its center of mass in the
plane of its elliptic orbit is described by the second
order ordinary quasilinear equation containing two
parameters, the eccentrity of orbit $e$, and the inertial
parameter $\mu$, which is small for the satellite being
almost dynamically symmetric [1].
We are interested in finding periodic solutions.
As $\mu=0$, the equation is explicitly integrable.
As $e=1$, it is singular.
A numerical study of the equation for $e$ close to $1$
is very difficult (see, e.g. [2]).

In the case of small $\mu$ the analysis of periodic
orbits is possible by applying the method of averaging,
or in the other words, computing the normal form
of the equation in the neigbourhood of periodic solutions
of the equation with $\mu=0$.
This approach has been proposed in [3] and developed
in [4]. Coefficients of the normal form are to be
found numerically as some integrals.

As $e\to 1$, there arises a problem in computing the
coefficients, because the straight algorithms of
computation are not regular in this case.
We propose a regular method of the computation,
based on integrating by a path in the complex plane and
factorizing integrands into functions
with close but diverse poles. Apart from the poles
we have a regular expression, which one can calculate
numerically.
Near the poles we use analytic regularization methods.

\begin{thebibliography}{4}
%
\bibitem{1}
V.V. Beletskii.
Motion of an Artificial Satellite about Its Center of Mass,
NASA TT F-429,
pp. 257--261.
%
\bibitem{2}
V.A. Sarychev.
Topics in satellite orientation.
Itogi nauki i tehniki, ser. "Issledovanije kosmicheskogo
prostranstva", Vol. 11. Moscow\-: VINITI\-, 1978.
(Russian)
%
\bibitem{3}
F.L. Chernous'ko.
Resonance Phenomena in the Motion of
a Satellite Relative to Its Mass Center.
USSR Computational Mathematics and Mathematical Phys., 3 (1964),
699-713.
%
\bibitem{4}
A.D. Bruno.
Local Methods in Nonlinear Differential Equations.
Springer-Verlag: Berlin-Heidelberg-New York-London-Paris-Tokyo,1989.
(Chapter~V).
%
\end{thebibliography}
\end{document}
